\documentclass[10pt,a4paper]{amsart}
\usepackage[margin=19mm]{geometry}
\usepackage{amsmath,amssymb,mathtools}
\usepackage[scr=rsfs,cal=euler]{mathalfa}
\usepackage{xfrac}
\usepackage[unicode,hidelinks,bookmarks=false]{hyperref}
\newcommand{\ie}{i.e.\ }
\newcommand{\cf}{cf.\ }
\newcommand{\resp}{respectively\ }
\newcommand{\Subsection}{\subsection}
\providecommand{\nobreakdash}{\nobreak\hbox{-}}
\setlength{\emergencystretch}{2em}
\multlinegap0pt
\usepackage{braket}
\usepackage{amssymb}
\usepackage{float}
\usepackage{xcolor}
\usepackage{algorithm}\usepackage{algpseudocode}
\usepackage{tikz}
\usepackage{bm}
\usepackage{mathtools}
\usepackage{dsfont}
\usepackage{csquotes}
\newtheorem{corollary}{Corollary}
\newtheorem{proposition}{Proposition}
\title{\textbf{Explicit M-Polynomial and Degree-Based Topological Indices of Generalized Hanoi Graphs}}
\author{El-Mehdi Mehiri}
\date{Source: arXiv:2511.12587v1 (CC BY 4.0)}
\begin{document}
\maketitle

\noindent\emph{Source and licence.} \textbf{Explicit M-Polynomial and Degree-Based Topological Indices of Generalized Hanoi Graphs}. Version arXiv:2511.12587v1 (CC BY 4.0), distributed by arXiv under the Creative Commons Attribution 4.0 International licence, http://creativecommons.org/licenses/by/4.0/, as recorded in the arXiv licence metadata for this exact version and shown on the abstract page https://arxiv.org/abs/2511.12587v1. Full licence notice: \texttt{licenses/arxiv-2511.12587-CC-BY-4.0.txt}.\par\smallskip

\noindent\textit{Editorial scope.} Counted unit: the article's proposition prop:eqdeg-top (source0.tex lines 1431--1452). The article's complete original proof of it is delivered with it (source0.tex lines 1454--1528, one \texttt{proof} environment of 75 lines). No mathematics is written, repaired or re-derived: the statement and the proof are the article's own lines, in the article's own notation, and no retained line is substituted or re-flowed.  Retained uncounted as the source-local context the proof needs: lines 1304--1308; lines 1376--1392.  Nothing was omitted from any retained span, so the package carries no omission record.  No retained line contains a citation, so the wrapper carries no bibliography and no bibliography entry.  No retained line contains a figure environment, an image-inclusion command or drawing code, so no image is delivered and the package declares no asset dependency.  Editorial formatting only, in addition: an AMS \texttt{amsart} preamble that restates the article's own declarations and macros used by the retained lines, namely the environments \texttt{corollary}, \texttt{proposition} on separate counters, together with the standard packages those lines need.  The retained environments print in this document as Corollary 1, Proposition 1 in order of appearance, whereas the article prints the same items with the numbering of its own full document.  The complete native source of the article as distributed in the arXiv e-print for this exact version is bundled unchanged as \texttt{sources/189208/source0.tex}.
\begin{align}
    |E_1(H_p^n)| &= \dfrac{1}{2} \sum_{\mu = 1}^{r} (p - \mu) \sum_{\nu = 0}^{\mu} \nu \, O_p^n(\nu \mid \mu),\label{eq:E1} \\
    |E_2(H_p^n)| &= \dfrac{1}{2} \binom{p}{2} p^n - \dfrac{1}{4} \sum_{\mu = 1}^{r} (p - \mu) \sum_{\nu = 0}^{\mu} \left(p + 3\mu - 2\nu - 1 \right) O_p^n(\nu \mid \mu),\label{eq:E2} \\
    |E_3(H_p^n)| &= \sum_{\mu = 1}^{r} (p - \mu) \sum_{\nu = 0}^{\mu} (\mu - \nu) \, O_p^n(\nu \mid \mu).\label{eq:E3}
\end{align}

\begin{corollary}\label{coro:another_espression_of_E2}  
An equivalent form of Equation~\eqref{eq:E2} is given by
\begin{equation}
    \begin{aligned}
|E_2(H_p^n)|&=\dfrac{1}{4} \sum_{\mu = 1}^{r}  \sum_{\nu = 0}^{\mu} \left( 2p(\nu- \mu)  + \mu( 3\mu- 2\nu - 1) \right) O_p^n(\nu \mid \mu)\\
&=\dfrac{1}{4} \sum_{\mu = 1}^{r}  \sum_{\nu = 0}^{\mu} \left( 2p\nu- 2p\mu  + 3\mu^2- 2\mu\nu - \mu \right) O_p^n(\nu \mid \mu).
\end{aligned}
\end{equation}
Moreover,
\begin{equation}
    \begin{aligned}
|E_1(H_p^n)|+ |E_2(H_p^n)|
&= \dfrac{1}{4} \sum_{\mu = 1}^{r}  \sum_{\nu = 0}^{\mu} \left( 3\mu^{2} -2p\mu-\mu+4p\nu-4\mu\nu \right) O_p^n(\nu \mid \mu)\\
&= \dfrac{1}{4} \sum_{\mu = 1}^{r}  \sum_{\nu = 0}^{\mu} \left( 4\nu(p-\mu)+\mu(3\mu-2p-1)  \right) O_p^n(\nu \mid \mu).
\end{aligned}
\end{equation}
\end{corollary}

\begin{proposition}\label{prop:eqdeg-top}  
For all $p \in \mathbb{N}_3$ and $n \in \mathbb{N}_0$,  let $r = \min\{n, p\}$ and  $\lambda=\Delta(H_p^n)=\binom{p}{2}$.
The number of edges $\{u,v\}\in E(H_p^n)$ satisfying $\deg(u)=\deg(v)=\lambda$ is
\begin{equation}\label{eq:sigma-top}
         \sigma(\lambda)=\begin{cases}
             \dfrac{1}{4}  \bigl(p^{2} -p\bigr)\, O_p^n(p), & \text{if } o(u)=o(v)=p,\\[4pt]
             \dfrac{1}{4}   \displaystyle\sum_{\nu = 0}^{p-1} 
             \bigl(p^{2}-5p+4\nu+4\bigr) O_p^n(\nu \mid p-1),
             & \text{if } o(u)=o(v)=p-1,\\[8pt]
             \displaystyle\sum_{\nu = 0}^{p-1}  (p-\nu-1)\, O_p^n(\nu \mid p-1),
             & \text{if } o(u)=p-1,\; o(v)=p.
         \end{cases}
\end{equation}
Consequently, we obtain the compact identity
\begin{equation}
        \sigma(\lambda)
        = \dfrac{1}{4}\,p
          \begin{Bmatrix}
             n+1\\p
          \end{Bmatrix}.
\end{equation}
\end{proposition}

\begin{proof}
If $\deg(u)=\deg(v)=\lambda=\Delta(H_p^n)$, then both endpoints have either all pegs occupied or all but one occupied. Hence,
\[
\{o(u),o(v)\}\in\bigl\{\{p,p\},\{p-1,p-1\},\{p-1,p\}\bigr\}.
\]
We analyze these cases separately.


\noindent
\begin{itemize}
    \item Case 1, $o(u)=o(v)=p$:

Here $\mu=f^{-1}(\lambda)=p$.  
Using Corollary~\ref{coro:another_espression_of_E2} with $\mu=p$, we find
\begin{align*}
\sigma(\lambda)
&= \dfrac{1}{4}   \sum_{\nu = 0}^{r} 
   \bigl(3p^{2} -2pp -p + 4p\nu -4p\nu\bigr) O_p^n(\nu \mid p)\\
&= \dfrac{1}{4}  (p^{2}-p)\,O_p^n(p).
\end{align*}

\item Case 2, $o(u)=o(v)=p-1$: 

Now $\mu=p-1$. Substituting $\mu=p-1$ in the same expression gives
\begin{align*}
\sigma(\lambda)
&= \dfrac{1}{4}   \sum_{\nu = 0}^{p-1} 
   \bigl(3(p-1)^2 -2p(p-1)-(p-1)+4p\nu-4\nu(p-1)\bigr)
   O_p^n(\nu \mid p-1)\\
&= \dfrac{1}{4}   \sum_{\nu = 0}^{p-1} 
   \bigl(p^{2}-5p+4\nu+4\bigr)\,O_p^n(\nu \mid p-1).
\end{align*}

\item  Case 3,  $\{o(u),o(v)\}=\{p-1,p\}$:

Here we use the expression from case~\texttt{(c\textsubscript{3})}, corresponding to a move from a multiton to an empty peg:
\begin{align*}
\sigma(\lambda)
&= \sum_{\nu = 0}^{p-1}(p-\mu)(\mu-\nu)\,O_p^n(\nu \mid \mu)
 = \sum_{\nu = 0}^{p-1}(p-(p-1))((p-1)-\nu)\,O_p^n(\nu \mid p-1)\\
&= \sum_{\nu = 0}^{p-1}(p-\nu-1)\,O_p^n(\nu \mid p-1).
\end{align*}
\end{itemize}
 
Summing the three contributions and applying the Stirling-type identity
\[
p\begin{Bmatrix} n\\p \end{Bmatrix}
   + \begin{Bmatrix} n\\p-1 \end{Bmatrix}
   = \begin{Bmatrix} n+1\\p \end{Bmatrix},
\]
we obtain the compact form

\begingroup
\allowdisplaybreaks
     \begin{align*}
         \sigma(\lambda)&= \dfrac{1}{4}  \left( p^{2} -p  \right) O_p^n(  p)+ \dfrac{1}{4}   \sum_{\nu = 0}^{p-1} \left( p^{2}-5p+4\nu+4 \right) O_p^n(\nu \mid p-1)+\sum_{\nu = 0}^{p-1}  (p-\nu-1) O_p^n(\nu \mid p-1)\\
         &=\dfrac{1}{4}  \left( p^{2} -p  \right) O_p^n(  p)+\dfrac{1}{4}  \left( p^{2} -p  \right)  \sum_{\nu = 0}^{p-1}   O_p^n(\nu \mid p-1)\\
         &=\dfrac{1}{4}  \left( p^{2} -p  \right) O_p^n(  p)+\dfrac{1}{4}  \left( p^{2} -p  \right)    O_p^n(  p-1)\\
         &=\dfrac{1}{4}  \left( p^{2} -p  \right) (O_p^n(  p)+ O_p^n(  p-1))\\
         &=\dfrac{1}{4}  \left( p^{2} -p  \right) \left( \begin{Bmatrix}
             n\\p
         \end{Bmatrix}p! +\begin{Bmatrix}
             n\\p-1
         \end{Bmatrix}(p-1)! \right)\\
         &=\dfrac{1}{4}  \left( p -1  \right) p! \left( p\begin{Bmatrix}
             n\\p
         \end{Bmatrix} +\begin{Bmatrix}
             n\\p-1
         \end{Bmatrix} \right)\\
         &=\dfrac{1}{4}  p\begin{Bmatrix}
             n+1\\p
         \end{Bmatrix}.\qedhere
     \end{align*}
\endgroup
\end{proof}

\end{document}
